\section{Discussion \& Limitations}
\label{sec:discussion}

\subsection{Calibration Sensitivity to Training Contamination}
Our contamination sweep (Section~\ref{sec:contamination}) reveals a critical boundary of our framework: while the PCMCI causal-graph backbone remains stable under sparse training poisoning, the empirical-CDF calibration step of H-CAS is vulnerable. As shown in Table~\ref{tab:contamination}, overall system performance degrades severely from 0.982 to 0.674 under just 1\% training contamination, dropping to 0.659 at 5\% contamination. In contrast, Isolation Forest (0.808) and PyTorch Autoencoder (0.843) degrade much more gracefully due to baseline outlier insulation. When attack events leak into the conformal calibration queue, the eCDF baseline is inflated, masking subsequent test attacks. Hardening the calibration step—using trimmed eCDFs or self-supervised anomaly filtering—is a critical open research direction.

\subsection{Scalability Limitations on Enterprise Features}
PCMCI causal discovery on the BETH Linux host telemetry dataset ($d = 882$ features, $T = 360$ training windows) requires 7,842 seconds (~2.18 hours) on a standard CPU. While this discovery phase is a one-time cost, this initialization cost scales quadratically with feature dimensions ($O(d^2 \tau_{\max} T)$). In high-density enterprise environments where unique system edges can reach thousands of features, this scaling delay represents a major limitation. Consequently, employing a variance-thresholding or mutual-information-based feature selection pre-pass is a mandatory requirement for deployment rather than an optional enhancement, ensuring only active system channels are monitored.

\subsection{Causal Sufficiency}
ZeroCausal assumes causal sufficiency---that no unobserved confounders jointly drive the observed provenance features. This is a standard working assumption in the PCMCI framework~\cite{pcmci2019} and in Pearl's structural causal hierarchy. In enterprise networks, unobserved external variables (e.g., domain controller group-policy pushes, NTP synchronization events, or user shift changes) could induce statistical correlations that PCMCI falsely learns as direct causal edges, resulting in localized false alarms.

However, ZeroCausal is \textit{robust to moderate violations} of this assumption for two reasons. First, the conformal prediction layer is distribution-free: it calibrates alarm thresholds empirically on the observed H-CAS score distribution, automatically absorbing any inflation caused by spurious causal edges. Second, the \texttt{AdaptiveWindowDetector} periodically refits the SCM, allowing stale or spurious edges to be discarded when the data distribution shifts. The practical consequence of confounded edges is elevated (but bounded) false alarms, not missed attacks---a favorable failure mode for a defense system.

\subsection{Linear Approximation of Count-Based Mechanisms}
The Linear SCM outperforms the Random Forest SCM on our datasets (Table~\ref{tab:performance}), which might appear to undermine the ``causal mechanism'' framing. We clarify: the \textit{causal graph} (which edges exist between provenance features) is the true invariant that ZeroCausal monitors, not the functional form $f_j$. The linear model serves as a first-order \textit{mechanism approximation}: on zero-inflated, sparse count data, parent influence is approximately proportional to parent event counts, and the residual distribution captures deviations. The Random Forest overfits rare count bursts in these heavy-tailed distributions, producing noisier residual rankings. On richer, continuous-valued event streams where non-linear interactions dominate, the Random Forest SCM may be preferable.

\subsection{Mimicry Attacks}
If an attacker is aware of the enterprise SCM, they could execute an APT that exactly mimics normal causal dependencies. However, successful mimicry requires the attacker to simultaneously satisfy \textit{all} $d$ causal equations across $\tau_{\max}$ time lags---a $d \times \tau_{\max}$-dimensional constraint. On OpTC ($d = 130$, $\tau_{\max} = 2$), this means matching 260 regression relationships while executing a multi-stage attack. This is a \textit{deliberate design trade-off}: ZeroCausal sacrifices defense against fully causal-aware adversaries in exchange for zero-label, online operation. In practice, real APT campaigns (lateral movement, privilege escalation, exfiltration) inject novel process--file edges that have no causal parents in the baseline, triggering structural novelty detection even without mechanism violations.

\subsection{Operational High-Recall Limitations on Real-world Data}
At a 95\% recall target, ZeroCausal exhibits a high false positive rate (FPR) of 93.40\% on the real enterprise OpTC logs. This means the system must flag 93.40\% of all normal windows just to catch 95\% of attacks, representing near-random performance and highlighting a critical operational limitation. 

To put this in perspective, under a standard Gaussian assumption for anomaly score distributions under $H_0$ and $H_1$, achieving an FPR below 50\% at a 95\% recall target requires a minimum classification AUROC of approximately 0.88 (a discrimination distance $d' \ge 1.645$). On real-world enterprise telemetry, which exhibits heavy-tailed count distributions and bursty administrative noise, the required AUROC is even higher. 

At this same operating point (FPR at 95\% recall), standard unsupervised baselines perform similarly poorly: Isolation Forest yields an FPR of 80.71\%, One-Class SVM has 75.17\%, and Local Outlier Factor (LOF) has 49.89\%. Only the PyTorch Autoencoder achieves a lower FPR of 6.65\% (at the cost of requiring clean training data and offline training). Causal-IDS, which utilizes fully labeled benign training data, achieves an FPR of 18.00\% (supervised reference) and 26.44\% (zero-label contaminated variant) on OpTC. 

This comparison underlines a fundamental property of anomaly-based IDS: catching every action of a stealthy, multi-stage attack inevitably flags standard benign administrative variability. Rather than optimizing for extreme recall targets at the cost of operator fatigue, ZeroCausal is designed to operate under a user-defined False Positive Rate (FPR) budget (e.g., 5\%). Conformal prediction bounds the empirical alarm rate by this budget ($4.09\% \pm 0.66\%$ across seeds), ensuring a constant, manageable alert volume. This demonstrates that conformal calibration provides \textit{stable empirical controls} even when ranking performance fluctuates.

\subsection{Ethical Considerations}
All real datasets (DARPA OpTC, StreamSpot, BETH) were collected from controlled testbeds or public repositories and do not contain private user information. TC3 and NODLINK are fully simulated. ZeroCausal is designed strictly for defensive monitoring.
